\section{Описание структуры и алгоритмов}
B-дерево хранит несколько упорядоченных записей в одном узле.
При минимальной степени $t\ge2$ некорневой узел содержит от $t-1$ до $2t-1$
ключей. Внутренний узел с $m$ ключами имеет $m+1$ детей; их диапазоны разделены
ключами родителя. Все листья находятся на одной глубине.
Корень имеет исключение по минимальному заполнению.

Поиск определяет позицию ключа и при отсутствии совпадения продолжается
в соответствующем ребёнке. При вставке сверху вниз заполненный ребёнок
разделяется до спуска: средняя запись поднимается в родителя, по $t-1$
записей остаётся в каждой половине.

При удалении перед спуском ребёнку обеспечивается запас до $t$ ключей
посредством заимствования через родителя либо слияния с соседом.
Для внутренней записи используются предшественник, преемник или слияние
минимальных детей. Пустой внутренний корень заменяется ребёнком.

Для непустого дерева высоты $h$ (число рёбер от корня до листа)
$n\ge2t^h-1$, откуда $h\le\log_t((n+1)/2)$.
При фиксированных $t$ и ограниченной длине слова основные операции
имеют логарифмическую сложность. Дополнительно учитывать стоимость
сравнений строк и перемещений внутри узла.
